\chapter{Общее операторное пространство и обменный мост дефектной пары}

\section{Постановка первой проверки Тома VI}

Том V установил два разных факта. Каждая ориентированная тёплицева ветвь
имеет индекс $\mp15$ и потому не допускает обратимого представителя в
своей фредгольмовой компоненте. Полная обменная Real-пара, напротив, имеет
нулевой обычный индекс, хотя представляет ненулевой класс пятнадцать в
$KO_6(M_{105}(\mathbb C)_{\mathbb R})$.

До начала вариационного анализа необходимо проверить, не были ли смешаны
три различных утверждения:
\begin{enumerate}
 \item невозможность обратить одну ориентированную ветвь;
 \item возможность обратить полный индекс-нулевой пакет;
 \item тривиальность или нетривиальность его Real/KО-класса.
\end{enumerate}

Это различение является обязательным: иначе можно либо ошибочно объявить
дефект неизбежным, либо, наоборот, принять обратимость за отсутствие
топологической материи.

\section{Сопоставление с проектом и литературой}

В Томе V были построены односторонний сдвиг $S$, его сопряжение $S^*$,
коэффициентный проектор $q_0$ ранга пятнадцать и обменная Real-форма.
Новая проверка не меняет эти данные и не вводит новый класс.

Общий факт теории Фредгольма гласит, что оператор индекса ноль допускает
конечноранговое возмущение до обратимого; см.
\path{arXiv:funct-an/9604005}. Унитарные дополнения с помощью дефектных
операторов являются стандартными матрицами Юлии--Халмоша; см.
\path{arXiv:2212.07808}. Компактное возмущение фредгольмова модуля не
обязано менять его $K$-гомологический класс; см.
\path{arXiv:1808.07911} и \path{arXiv:hep-th/0012046}. Следовательно,
обратимый представитель и нулевой класс не являются синонимами.

Отличие от закрытых конструкций проекта состоит в разрешении одного
конкретного конечномерного канала между сопряжёнными ориентациями. Никакой
такой канал не использовался в ориентированных индексных нижних границах.

\section{Три пространства конфигураций}

Обозначим через $\mathcal H_q=H^2(S^1)\otimes q_0\mathbb C^{105}$
коэффициентный модуль ранга пятнадцать и положим
\begin{equation}
 S_q=S\otimes I_{15},
 \qquad
 P_q=|e_0\rangle\langle e_0|\otimes I_{15}.
\end{equation}

Нужно различать:
\begin{description}
 \item[$\mathfrak C_{\rm or}$] блок-диагональное пространство, в котором
 две ориентации сохраняются по отдельности;
 \item[$\mathfrak C_{\rm ex}$] обменное пространство, допускающее
 конечномерные межветвевые операторы;
 \item[$\mathfrak C_{0}$] пространство физического пустого вакуума,
 где должны быть дополнительно фиксированы представление алгебры,
 Real/Clifford-структура и нулевой класс.
\end{description}

В $\mathfrak C_{\rm or}$ оператор
\begin{equation}
 T_0=
 \begin{pmatrix}
  S_q&0\\
  0&S_q^*
 \end{pmatrix}
 \label{eq:v6-balanced-pair-t0}
\end{equation}
имеет
\begin{equation}
 \dim\ker T_0=15,
 \qquad
 \dim\operatorname{coker}T_0=15,
 \qquad
 \operatorname{ind}T_0=0.
\end{equation}
Каждый диагональный блок сохраняет индекс $\mp15$, поэтому
ориентационно-блочное возмущение не может сделать оба блока обратимыми.

\section{Явный обменный мост}

Разрешим конечномерный оператор
\begin{equation}
 K=
 \begin{pmatrix}
  0&P_q\\
  0&0
 \end{pmatrix},
 \qquad
 \operatorname{rank}K=15,
 \end{equation}
который отображает ядро ветви $S_q^*$ на коядро ветви $S_q$. Рассмотрим
однопараметрическое семейство
\begin{equation}
 T_\lambda=T_0+\lambda K
 =
 \begin{pmatrix}
  S_q&\lambda P_q\\
  0&S_q^*
 \end{pmatrix},
 \qquad 0\leq\lambda\leq1.
 \label{eq:v6-exchange-bridge-path}
\end{equation}

Используя
\begin{equation}
 S_q^*S_q=I,
 \qquad
 S_qS_q^*=I-P_q,
 \qquad
 S_q^*P_q=P_qS_q=0,
\end{equation}
получаем
\begin{align}
 T_\lambda^*T_\lambda
 &=
 \begin{pmatrix}
 I&0\\0&I+(\lambda^2-1)P_q
 \end{pmatrix},\\
 T_\lambda T_\lambda^*
 &=
 \begin{pmatrix}
 I+(\lambda^2-1)P_q&0\\0&I
 \end{pmatrix}.
 \label{eq:v6-exchange-bridge-products}
\end{align}

\begin{theorem}[Обратимое дополнение нейтральной пары]
Для каждого $\lambda>0$ оператор $T_\lambda$ обратим, его минимальное
сингулярное число равно $\min(1,\lambda)$, а при $\lambda=1$ оператор
$T_1$ унитарен. Следовательно, ядро и коядро сбалансированной пары могут
быть удалены конечноранговым межветвевым мостом без выхода из пространства
фредгольмовых операторов.
\end{theorem}

В конечномерной прямоугольной модели этот оператор является точной
перестановочной унитарной матрицей. После канонического обменного
отождествления двух ветвей она становится самосопряжённой унитарной
матрицей. Машинный аудит проверяет эти тождества для коэффициентного ранга
пятнадцать.

\section{Почему класс пятнадцать при этом не исчезает}

Возмущение $K$ компактно. В комплексной фредгольмовой категории оно
устраняет явные нулевые моды конкретного представителя, но не обязано
превращать фредгольмов модуль в вырожденный нулевой цикл. В частности,
из одной обратимости не следует изменение:
\begin{enumerate}
 \item представление коэффициентной алгебры;
 \item символ $V_{15}$ и его Real-условие;
 \item Clifford-степень;
 \item класс комплексification $(-15,+15)$.
\end{enumerate}

Поэтому доказанный в Томе V Real-класс пятнадцать в принципе совместим с
обратимым представителем полного пакета: ненулевой $K$-класс измеряет не
только размер ядра конкретного оператора. Однако утверждать, что именно
$T_\lambda$ сохраняет проектный Real-класс, можно только после проверки
полного $J$ и Clifford-действия. Настоящая проверка доказывает комплексное
конечноранговое дополнение и обменную самосопряжённость конечной модели,
но не подменяет ими следующий Real-аудит.

\begin{proposition}[Уточнение статуса дефицита $1/7$]
Нижняя граница $1/7$ остаётся строгой для каждой ориентированной ветви и
для блок-диагонального пространства $\mathfrak C_{\rm or}$. Она не является
нижней границей суммы ядра и коядра во всём расширенном пространстве
$\mathfrak C_{\rm ex}$: мост $K$ делает эту сумму равной нулю. Сохранение
числа $1/7$ в полном пакете относится к нормированному Real/KО-классу или
к выбранному блочному представителю, а не к неизбежной необратимости
любого представителя пары.
\end{proposition}

Это уточнение не отменяет класс пятнадцать. Оно запрещает отождествлять
три разных понятия: топологический класс, явные нулевые моды и физическую
материю.

\section{Два возможных механизма рождения}

Полученный результат разделяет дальнейшую программу на две ветви.

\subsection*{Механизм A: создание класса}

Если физический вакуум является вырожденным нулевым $K$-циклом, переход к
классу пятнадцать должен менять Real/KО-класс. Тогда необходим выход из
пространства допустимых символов, потеря фредгольмовости или изменение
граничной задачи. Прежний сценарий сингулярного рождения относится именно
к этой ветви.

\subsection*{Механизм B: проявление скрытого класса}

Если родительский вакуум уже несёт обратимый представитель класса
пятнадцать, то локальные дефекты могут проявиться при выключении
межветвевого моста:
\begin{equation}
 T_1\longrightarrow T_\lambda\longrightarrow T_0,
 \qquad \lambda\downarrow0.
 \end{equation}
В этом механизме общий класс не рождается из нуля; меняется способ его
локального проявления. Он ранее не был проверен проектом и не должен быть
отброшен до анализа допустимых физических одноформ.

\section{Что эта проверка не доказывает}

Пока не установлено, что $K$ принадлежит физическому бимодулю одноформ,
сохраняет полный проектный $J$, калибровочную группу, BV/BRST-фактор и
бюджет $H_{15}$. Не выведен также потенциал для параметра $\lambda$.
Следовательно, найдено общее операторное пространство, но ещё не физическая
динамика.

\begin{nogocriterion}[Запрет преждевременного выбора механизма]
Нельзя требовать нефредгольмова рождения класса, пока не проверено, что
родительский вакуум имеет нулевой Real/KО-класс. Нельзя также объявлять
механизм скрытого класса физическим, пока мост $K$ не получен из
допустимой одноформы и одного родительского функционала.
\end{nogocriterion}

\begin{protocol}
\begin{enumerate}
 \item индекс отдельной ориентации: \textbf{$\mp15$};
 \item обычный индекс полной пары: \textbf{ноль};
 \item блок-диагональная обратимость: \textbf{невозможна};
 \item ранг минимального межветвевого моста: \textbf{15};
 \item обратимость $T_\lambda$ при $\lambda>0$: \textbf{доказана};
 \item унитарность $T_1$: \textbf{доказана};
 \item выход из фредгольмова пространства для устранения нулевых мод пары:
 \textbf{не требуется};
 \item обнуление класса из одной обратимости: \textbf{не следует};
 \item сохранение проектного Real/KО-класса конкретным мостом:
 \textbf{требует следующего проверки};
 \item дефицит $1/7$ как всеобщая граница ядра и коядра полной пары:
 \textbf{отозван};
 \item дефицит $1/7$ для ориентированной ветви: \textbf{сохранён};
 \item общее операторное пространство: \textbf{построено};
 \item физическая допустимость моста: \textbf{не проверена};
 \item следующая проверка:
 \textbf{родительская допустимость обменного моста}.
\end{enumerate}
\end{protocol}

Машинный сертификат сохранён в
\path{s2t_v6_common_configuration_space_gate_results.json}.